

% ---------------------------------------------------------
% Cores básicas e fundo de página
% ---------------------------------------------------------
\usepackage[xetex,dvipsnames,svgnames,table]{xcolor}

% Cores definidas por CMYK
% \definecolor{fondpaille}{cmyk}{0,0.02,0.2,0}
\definecolor{fondpaille}{cmyk}{0, 0.0118, 0.1176, 0, 0.01} %cmyk(0%, 2.75%, 23.53%, 0%, 1%)
\definecolor{ferrugem}{cmyk}{0, 0.4457, 0.75, 0.6392}

% Cores definidas por hexadecimal
\definecolor{moleskin}{HTML}{FFF8DC}
\definecolor{laranja}{HTML}{F07241}
\definecolor{marron}{HTML}{800000}
\definecolor{lovely}{HTML}{C04848}
\definecolor{click}{HTML}{480048}
\definecolor{tinta-velha}{HTML}{4e4643}
\definecolor{Browned_Cherry}{HTML}{910831}
\definecolor{azul}{HTML}{000080}

% Cores da capa
\definecolor{BLACK}{HTML}{000000}
\definecolor{ELECTRICRED}{HTML}{DD0000}
\definecolor{TANGERINYELLOW}{HTML}{FFCE00}

% Cor consolidada para texto principal
\colorlet{cor-da-tinta}{ferrugem}
\newif\ifferrugem
\ferrugemfalse

% Cor de fundo da página
\pagecolor{fondpaille}

% ---------------------------------------------------------
% Fontes
% ---------------------------------------------------------
\usepackage{fontspec} % XeLaTeX/LuaLaTeX fonts
\setmainfont{EBGaramond}[ % Tudo junto, não 'EB Garamond'
  %% TECTONIC:  %%  %%  %%  %%  %%
  Path=./auxiliares/ebgaramond/,
  Extension=.ttf,
  UprightFont=*-Regular,
  BoldFont=*-Bold,
  ItalicFont=*-Italic,
  BoldItalicFont=*-BoldItalic,
  %%  %%  %%  %%  %%  %%  %%  %%  %%
  Numbers=OldStyle,
  RawFeature={
    -aalt,
    -c2pc,
    -c2sc,
    -case,
    +dlig,
    -frac,
    -hist,
    +hlig,
    +kern,
    +liga,
    +mark,
    +mkmk,
    +onum,
    +pnum,
    +ss05,
    +cv06
  }
]
\usepackage[math-style=ISO, bold-style=ISO]{unicode-math}
% \setmathfont{OldStandard-Math.otf}
% \setmathfont{Garamond-Math.otf}[StylisticSet={5,6,7,9}]
\setmathfont{NewCMMath-Book.otf}

% ---------------------------------------------------------
% Línguas e hifenização
% ---------------------------------------------------------
\usepackage{polyglossia}
\setdefaultlanguage{portuguese}
\PolyglossiaSetup{portuguese}{indentfirst=true}
\setotherlanguages{english,french} % TECTONIC: 'german' não funciona

% ---------------------------------------------------------
% Tipografia avançada
% ---------------------------------------------------------
\usepackage{microtype}      % ajustes de tipografia
\usepackage[bottom]{footmisc} % pés de página alinhados ao fundo
\linespread{1.3}            % espaçamento de linha

% ---------------------------------------------------------
% Matemática e símbolos
% ---------------------------------------------------------
\usepackage{amsmath,dsfont,mathrsfs,wasysym}
\usepackage[makeroom,thicklines]{cancel} % para riscar
\usepackage{esvect}       % vetores

% ---------------------------------------------------------
% Código-fonte e Python
% ---------------------------------------------------------
\usepackage{fvextra}
\usepackage{fancyvrb}
\usepackage{listings}
% \usepackage{pythontex}  TECTONIC: Não é compatível

% ---------------------------------------------------------
% Gráficos e figuras
% ---------------------------------------------------------
\usepackage{graphicx}
\usepackage{svg}
\usepackage{tikz}
\usepackage{tkz-fct}
\usepackage{tikzpagenodes}
\usepackage{circuitikz}
\usepackage{pgfplots,pgfplotstable}
\pgfplotsset{compat=newest}
\usetikzlibrary{
  calc,
  positioning,
  automata,
  arrows,
  shadows,
  shapes.multipart,
  matrix,
  backgrounds,
  decorations.pathreplacing,
  fit,
  plotmarks,
  tikzmark
}
\pgfdeclarelayer{background}
\pgfdeclarelayer{foreground}
\pgfsetlayers{background,main,foreground}

% ---------------------------------------------------------
% Tabelas avançadas
% ---------------------------------------------------------
\usepackage{array,longtable,booktabs,bigstrut,tabularx,seqsplit}

% ---------------------------------------------------------
% Notas de rodapé e citações
% ---------------------------------------------------------
\usepackage{csquotes}
\usepackage[toc,page]{appendix}
\renewcommand\appendixpagename{Apêndices}
\usepackage[
  backend=biber,
  style=numeric-comp,
  language=auto,
  natbib=false,
  block=space,
  isbn=false,
  url=true,
  doi=false,
  mcite=true,
  eprint=false,
  sorting=none
]{biblatex}
\addbibresource{bibliografia.bib}

% Formato das notas de rodapé
\renewcommand{\thefootnote}{\alph{footnote}}
\makeatletter
\renewcommand\@makefnmark{\@textsuperscript{\normalfont(\@thefnmark)}}
\makeatother
\long\def\symbolfootnote[#1]#2{%
  \begingroup
  \def\thefootnote{\fnsymbol{footnote}}
  \footnote[#1]{#2}
  \endgroup
}

% ---------------------------------------------------------
% Secções e sumário
% ---------------------------------------------------------
\usepackage{titlesec}
\usepackage{tocloft}
\usepackage{etoolbox}

% Novos níveis de secção
\titleclass{\subsubsubsection}{straight}[\subsubsection]
\titleclass{\subsubsubsubsection}{straight}[\subsubsubsection]
\newcounter{subsubsubsection}[subsubsection]
\newcounter{subsubsubsubsection}[subsubsubsection]
\renewcommand{\thesubsubsubsection}{\thesubsubsection.\arabic{subsubsubsection}}
\renewcommand{\thesubsubsubsubsection}{\thesubsubsubsection.\arabic{subsubsubsubsection}}

% Formato visual
\titleformat{\subsubsubsection}{\normalfont\normalsize\bfseries}{\thesubsubsubsection}{1em}{}
\titlespacing*{\subsubsubsection}{0pt}{3.25ex plus 1ex minus .2ex}{1.5ex plus .2ex}
\titleformat{\subsubsubsubsection}{\normalfont\normalsize\bfseries}{\thesubsubsubsubsection}{1em}{}
\titlespacing*{\subsubsubsubsection}{0pt}{3.25ex plus 1ex minus .2ex}{1.5ex plus .2ex}

% Sumário detalhado
\renewcommand{\cftsecfont}{\bfseries}
\renewcommand{\cftsubsecfont}{\normalfont}
\renewcommand{\cftsubsubsecfont}{\normalfont}
\makeatletter
\renewcommand{\l@section}{\@dottedtocline{1}{1.5em}{1.5em}}
\renewcommand{\l@subsection}{\@dottedtocline{2}{2.5em}{2em}}
\renewcommand{\l@subsubsection}{\@dottedtocline{3}{4em}{2.5em}}
\newcommand{\l@subsubsubsection}{\@dottedtocline{4}{6em}{3em}}
\newcommand{\l@subsubsubsubsection}{\@dottedtocline{5}{8em}{4em}}
\providecommand*{\toclevel@subsubsubsection}{4}
\providecommand*{\toclevel@subsubsubsubsection}{5}
\makeatother
\setcounter{secnumdepth}{5}
\setcounter{tocdepth}{5}

% ---------------------------------------------------------
% Layout de página
% ---------------------------------------------------------
\usepackage[a4paper,left=2.5cm,right=2.5cm,top=2.5cm,bottom=2.5cm,includehead,includefoot]{geometry}
\usepackage{setspace} % espaçamento de linhas
\usepackage{lettrine,calligra} % letras capitulares
\renewcommand{\LettrineFontHook}{\calligra}
\LettrineTextFont{\itshape}
\setcounter{DefaultLines}{3}

% Prevenção de viúvas e órfãos
\clubpenalty=100
\linepenalty=100
\widowpenalty=100

% Outros ajustes tipográficos
\usepackage{xurl}
\usepackage[hang,small,bf]{caption}
\setlength{\captionmargin}{1.5in}
\renewcommand{\fboxrule}{2pt}

% ---------------------------------------------------------
% Utilitários diversos
% ---------------------------------------------------------
\usepackage{comment}
\usepackage{enumerate}
\newcounter{lastenum}
\newcommand{\storelast}{\setcounter{lastenum}{\value{enumi}}}
\newcommand{\uselast}{\setcounter{enumi}{\value{lastenum}}}
\newcommand\CC{C\nolinebreak[4]\hspace{-.05em}\raisebox{.4ex}{\relsize{-3}{\textbf{++}}}}
\usepackage{rotating}
\usepackage{array}
\usepackage{lastpage}
\usepackage{refcount}
\usepackage{standalone}
\usepackage{relsize}
\usepackage[normalem]{ulem}
\usepackage[most]{tcolorbox}

% ---------------------------------------------------------
% Datas e formatação numérica
% ---------------------------------------------------------
\usepackage[nodayofweek,level]{datetime}
\usepackage[output-decimal-marker={,}]{siunitx}

% ---------------------------------------------------------
% Pacotes de cálculo e exercícios
% ---------------------------------------------------------
\usepackage{xlop}      % cálculo automático

% ---------------------------------------------------------
% Topo e rodapé de páginas (logos e texto)
% ---------------------------------------------------------
\usepackage{eso-pic}

\AddToShipoutPictureBG{%           % Início do bloco que adiciona elementos ao fundo da página
  \ifnum\value{page}>2           % Condição: só aplica a partir da página 3
    \AtPageUpperLeft{%            % Posiciona tudo relativamente ao canto superior esquerdo da página
      \begin{tikzpicture}[remember picture, overlay]% Início do ambiente TikZ sobreposto (overlay) e lembrando posição
        % Linhas superior e inferior
        \draw[line width=0.5pt, color=black] (2.5cm, -3cm) -- (\paperwidth - 2.5cm, -3cm); % Linha horizontal superior
        \draw[line width=0.5pt, color=black] (2.5cm, -\paperheight + 3cm) -- (\paperwidth - 2.5cm, -\paperheight + 3cm); % Linha horizontal inferior
        
        % Conteúdo página ímpar
        \ifodd\value{page}           % Verifica se a página é ímpar
          \node [anchor=west, text depth=0pt] at (2.5cm, -2.7cm){\tiny\href{mailto:luisav.ferreira@gmail.com}{Grupo Disciplinar 510}}; % Texto canto superior esquerdo
          \node [anchor=north east] at (\paperwidth-2.5cm, -2.2cm){\includegraphics[width=0.05\paperwidth]{logos/aeB05.pdf}}; % Logo canto superior direito
          \node [anchor=west, text height=0pt] at (2.5cm, -\paperheight + 2.7cm){\tiny\href{mailto:luisav.ferreira@gmail.com}{Luís A.V. Ferreira}}; % Texto canto inferior esquerdo
          \node [anchor=west] at (2.5cm, -\paperheight + 2.5cm){\tiny\jobname.tex \smash{\rule[-0.5ex]{0.6pt}{1ex}}}; % Nome do ficheiro e linha decorativa
        % Conteúdo página par
        \else                        % Caso seja página par
          \node [anchor=west, text depth=0pt] at (2.5cm, -2.7cm){\includegraphics[width=0.05\paperwidth]{logos/aeB05.pdf}}; % Logo canto superior esquerdo
          \node [anchor=east, text depth=0pt] at (\paperwidth - 2.5cm, -2.7cm){\tiny\href{mailto:luisav.ferreira@gmail.com}{Grupo Disciplinar 510}}; % Texto canto superior direito
          \node [anchor=east, text height=0pt] at (\paperwidth - 2.5cm, -\paperheight + 2.7cm){\tiny\href{mailto:luisav.ferreira@gmail.com}{Luís A.V. Ferreira}}; % Texto canto inferior direito
          \node [anchor=east] at (\paperwidth - 2.5cm, -\paperheight + 2.5cm){\tiny\jobname.tex \smash{\rule[-0.5ex]{0.6pt}{1ex}}}; % Nome do ficheiro e linha decorativa
        \fi
      \end{tikzpicture}
    }
  \fi
}

% ---------------------------------------------------------
% hyperref
% ---------------------------------------------------------

\usepackage[%
  bookmarks=true,
  citecolor=black,
  filecolor=black,
  linkcolor=black,
  urlcolor=blue,     % Cor específica para URLs
  unicode=false,
  colorlinks=false,
  plainpages=false,
  pdfpagelabels,
  colorlinks=true,  % Ativa cores em vez de caixas
  xetex]{hyperref} 
